\documentclass[10pt,a4paper]{amsart}
\usepackage[margin=19mm]{geometry}
\usepackage{amsmath,amssymb,mathtools}
\usepackage[scr=rsfs,cal=euler]{mathalfa}
\usepackage{xfrac}
\usepackage[unicode,hidelinks,bookmarks=false]{hyperref}
\newcommand{\ie}{i.e.\ }
\newcommand{\cf}{cf.\ }
\newcommand{\resp}{respectively\ }
\newcommand{\Subsection}{\subsection}
\providecommand{\nobreakdash}{\nobreak\hbox{-}}
\setlength{\emergencystretch}{2em}
\multlinegap0pt
\usepackage{bm}
\usepackage{bbm}
\usepackage{dsfont}
\usepackage{xspace}
\usepackage{cancel}
\usepackage{mathtools}
\usepackage{amsthm}
\makeatletter
\@ifundefined{theorem}{\newtheorem{theorem}{Theorem}}{}
\@ifundefined{lemma}{\newtheorem{lemma}[theorem]{Lemma}}{}
\@ifundefined{proposition}{\newtheorem{proposition}[theorem]{Proposition}}{}
\makeatother
\DeclareMathOperator{\mes}{mes}
\newcommand{\balp}{\boldsymbol{\alpha}}
\newcommand{\cM}{{\mathcal M}}
\title[On the density of rational lines on diagonal cubic hypersurf]{On the density of rational lines on diagonal cubic hypersurfaces, II}
\author{Scott Parsell, Kiseok Yeon}
\date{Source: arXiv:2605.30549v2 (CC BY 4.0)}
\begin{document}
\maketitle

\noindent\emph{Source and licence.} On the density of rational lines on diagonal cubic hypersurfaces, II. Version arXiv:2605.30549v2 (CC BY 4.0), distributed under the Creative Commons Attribution 4.0 International licence, as recorded in the arXiv licence metadata for this exact version and shown on the abstract page https://arxiv.org/abs/2605.30549v2. Full rights notice: licenses/arxiv-2605.30549-CC-BY-4.0.txt.\par\smallskip

\noindent\textit{Editorial scope.} Counted unit: the Lemma whose source label is \texttt{prop3.2} in the article (source0.tex lines 1108--1114, source label \texttt{prop3.2}), delivered together with the article's own complete original proof of it (source0.tex lines 1115--1312, one \texttt{proof} environment of 198 lines). No mathematics is written, repaired or re-derived: statement and proof are the article's own text line for line. Required context retained uncounted: J10 (lines 463--465); basic minor (lines 524--529); def of Ll(L;H) (lines 603--612); pro5.2 (lines 620--630); improved minor (lines 688--696); k-1 dim measure (lines 1000--1006). Every definition, display and label of those lines is retained unchanged. Omitted from this package and not redistributed: 1--462, 466--523, 530--602, 613--619, 631--687, 697--999, 1007--1107, 1313--1752. Nothing inside a retained excerpt is omitted or altered: every retained body is delivered exactly as the article's lines are written, so no omission receipt of any kind accompanies this package. Citations: the retained excerpts carry no citation command at all, so no bibliography entry of the article is transcribed and no reference is redistributed. Reference adaptation: none was needed and none was made; every reference inside the delivered unit resolves inside it, so no unresolved reference or citation is produced. Printed slips kept literal: no printed word, symbol, hypothesis or number of the counted statement or of its proof was altered. Layout and class: the article's own class is \texttt{amsart} with packages including tikz, xcolor, amssymb, latexsym, amsmath, extarrows, graphicx, mathrsfs, comment, hyperref, url, amstext, bbm; the wrapper is the lane's pinned \texttt{amsart} editorial wrapper at 10pt on A4, which restates the theorem-like environments the retained text needs and adds the article's own macro definitions that the retained text needs (\texttt{\textbackslash balp}, \texttt{\textbackslash cM}, \texttt{\textbackslash mes}), with extra wrapper packages (bm, bbm, dsfont, xspace, cancel, mathtools). The delivered numbering is the unit's own. Assets: no retained span contains a figure environment, a graphics-inclusion command, a PDF-page inclusion, a table or a TikZ picture; no image file is copied into this package, the package declares no asset dependency and no image file is redistributed with it. Licence: version arXiv:2605.30549v2 of this article is distributed under the Creative Commons Attribution 4.0 International licence, as recorded in the arXiv licence metadata for this exact version and shown on the abstract page https://arxiv.org/abs/2605.30549v2. A full rights notice is delivered at licenses/arxiv-2605.30549-CC-BY-4.0.txt. Characters: every retained excerpt is pure ASCII, so the delivered engine, XeLaTeX with its default Latin Modern text font, reports no missing character.
\begin{equation}\label{J10}
J_{10}(X) \ll X^{20+\varepsilon}.
\end{equation}

\begin{lemma}\label{basic minor}
    Let $Q$ and $X$ be positive numbers with $Q \le X^{3/2}$, and suppose that $s$ is a natural number. Then for $1 \le l \le 4$ one has
    \begin{equation*}\label{meanvalue lemma 2.2}
        \int_{\cM_l(Q)} |F(\boldsymbol{\alpha})|^{2s}\, d\boldsymbol{\alpha} \ll X^{8} \cdot J_s(2X) \cdot (Q^{-1} + X^{-1})(\log X)^{4s+1}. 
     \end{equation*}
\end{lemma}

\begin{equation}\label{def of Ll(L;H)}
    \begin{aligned}
        \mathcal{L}_l(L;H) &=\left\{\begin{aligned}
            &\mathcal{L}_{2,3}^*(L; H)  \ \ \ \text{when}\ l=1, 3\\
            &\mathcal{L}_{3,2}^*(L; H)\ \ \ \text{when}\ l=2, 4.\\
          %  &(2,3), \ \ \ \text{when}\ l=3\\
          %  &(3,2),\ \ \ \text{when}\ l=4
        \end{aligned}\right.
    \end{aligned}
\end{equation}

\begin{proposition}\label{pro5.2}
Let $H,L$ and $X$ be positive numbers with $H\leq X^{3/2}$ and $L\leq H^{1/2}$. 
  Consider the set
  \begin{equation*}
  \mathfrak{L}(L;H) :=  \{(\alpha,\beta)\in [0,1)^2:\ \alpha\in \mathfrak{M}(H),\ \beta\in \mathfrak{L}^{(\alpha)}(L;H)\}.
  \end{equation*}
  Then, the set  $\mathfrak{L}(L;H)$ is measurable and 
  \begin{equation*}
      \mes{\left(  \mathfrak{L}(L;H)\right)}\ll   L^2H^3X^{-6}.
  \end{equation*}
\end{proposition}

\begin{lemma}\label{improved minor}
    Let $Q, H, L$, and $X$ be positive numbers with $Q \leq H\leq X^{3/2}$ and $L\leq H^{1/2}$, and suppose that $s $ is a natural number. Then for $1\leq l\leq 4$ one has 
    \begin{equation*}\label{sharp bound for T}
    \begin{aligned}
        & \int_{\mathcal{M}_l(Q)\cap \mathcal{L}_{l}(L;H)} |F(\boldsymbol{\alpha})|^{2s}\, d\boldsymbol{\alpha}  \\
        & \ \ \ \ \ \ \ \ll X^{8} \cdot J_s(2X)\cdot \left(Q^{-1}+X^{-1}\right)\bigl(L^{-1}+X^{-1/2}\bigr)(\log X)^{4s+2}.
    \end{aligned}
    \end{equation*}
\end{lemma}

\begin{lemma}\label{k-1 dim measure}
    Suppose that $\mathfrak{D}$ is a measurable set in $[0,1)^{2}$.  Then one has
    \begin{equation*}
        \int_0^1\int_{\mathfrak{D}}\int_0^1 |F(\boldsymbol{\alpha})|^{8} \, d\boldsymbol{\alpha}\ll X^{10+\varepsilon}\cdot \mes{(\mathfrak{D})},
    \end{equation*}
where $d\boldsymbol{\alpha}=d\alpha_1 \, d\alpha_2 \, d\alpha_3 \, d\alpha_4.$
\end{lemma}

\begin{lemma}\label{prop3.2}
    One has
    \begin{equation*}\label{claim in lemma 3.2}
\int_{\mathfrak{n}_{\delta}}|F(\boldsymbol{\alpha})|^{18}d\boldsymbol{\alpha}\ll X^{24-\nu},
    \end{equation*}
    for some $\nu>0,$ where $d\boldsymbol{\alpha}=d\alpha_1 \, d\alpha_2 \, d\alpha_3 \, d\alpha_4.$
\end{lemma}

\begin{proof}
 Recall the definition of $\mathfrak{M}(H)$. We observe that for any $H\leq X^{\delta/100},$ one has
\begin{equation}\label{inclusion}
    \mathfrak{M}(H)^4\subseteq \mathfrak{N}_{\delta}.
\end{equation}
Observe that for $H>0,$ one has
\begin{equation}\label{pruning argument crucial}
\begin{aligned}
   \mathfrak{M}(H)^4 \setminus \mathfrak{M}(H/2)^4
   =\mathfrak{P}_1(H)\cup \mathfrak{P}_2(H)\cup\mathfrak{P}_3(H)\cup \mathfrak{P}_4(H),
\end{aligned}
\end{equation}
where 
\begin{equation*}
    \begin{aligned}
        \mathfrak{P}_1(H)&=\mathfrak{M}(H)^4\setminus(\mathfrak{M}(H/2)\times\mathfrak{M}(H)^3)\\
 \mathfrak{P}_2(H)&=  (\mathfrak{M}(H/2)\times\mathfrak{M}(H)^3)\setminus (\mathfrak{M}(H/2)^2\times\mathfrak{M}(H)^2)\\
  \mathfrak{P}_3(H)&= (\mathfrak{M}(H/2)^2\times\mathfrak{M}(H)^2)\setminus (\mathfrak{M}(H/2)^3\times\mathfrak{M}(H))\\
  \mathfrak{P}_4(H)&=  (\mathfrak{M}(H/2)^3\times\mathfrak{M}(H))\setminus \mathfrak{M}(H/2)^4.\end{aligned}
\end{equation*}
%To verify the inequality $(\ref{claim in lemma 3.2})$, 
It suffices to show that for $1\leq l\leq 4$ and for all $X^{\delta/100}\leq H\leq X^{3/2}$ with $\delta=10^{-10}$, one has
\begin{equation}\label{claim}
\int_{ \mathfrak{P}_l(H)}|F(\boldsymbol{\alpha})|^{18} \, d\boldsymbol{\alpha}\ll X^{24-\gamma}\ \text{for some}\ \gamma>0.
\end{equation}
In fact, note that by (\ref{inclusion}) one has
\begin{equation}\label{fdsa}
\begin{aligned}
    \mathfrak{n}_{\delta}&\subseteq \mathfrak{M}(X^{3/2})^4\setminus \mathfrak{M}(X^{\delta/100})^4\\
    &\subseteq\bigcup_{j=0}^{\rm L}(\mathfrak{M}(2^{-j}X^{3/2})^4\setminus \mathfrak{M}(2^{-j-1}X^{3/2})^4), 
\end{aligned}
\end{equation}
for some ${\rm L}=O(\log X).$
Then, it follows by $(\ref{pruning argument crucial})$ that 
\begin{equation}\label{asdf1}
\begin{aligned}
    \mathfrak{n}_{\delta}&\subseteq \bigcup_{j=0}^{\rm L}\bigl(\mathfrak{P}_1(2^{-j}X^{3/2})\cup \mathfrak{P}_2(2^{-j}X^{3/2})\cup\mathfrak{P}_3(2^{-j}X^{3/2})\cup \mathfrak{P}_4(2^{-j}X^{3/2})\bigl).
\end{aligned}
\end{equation}
Then, assuming the truth of $(\ref{claim})$, we infer from $(\ref{asdf1})$  that 
\begin{equation*}\label{reasoning}
\begin{aligned}
\int_{\mathfrak{n}_{\delta}}|F(\boldsymbol{\alpha})|^{18}\, d\boldsymbol{\alpha}&\ll \sum_{j=0}^{\rm L}\sum_{l=1}^4\int_{ \mathfrak{P}_l(2^{-j}X^{3/2})}|F(\boldsymbol{\alpha})|^{18} \, d\boldsymbol{\alpha} \ll X^{24-\nu},
\end{aligned}
\end{equation*}
for some $\nu > 0$, as required.
Hence, we turn to verifying $(\ref{claim}).$
Recall the definition (\ref{def of Ll(L;H)}) of $ \mathcal{L}_{l}(L;H)$ with $1\leq l\leq 4.$ We denote $[0,1)^4\setminus \mathcal{L}_{l}(L;H)$ by $(\mathcal{L}_{l}(L;H))^c$, and we set $L=H^{1/7}$. 
Then, we see that
\begin{equation}\label{new dissection}
\begin{aligned}
&\int_{\mathfrak{P}_l(H)}|F(\boldsymbol{\alpha})|^{18} \, d\balp
\\
&=\int_{\mathfrak{P}_l(H)\cap \mathcal{L}_{l}(L;H)}|F(\boldsymbol{\alpha})|^{18}\, d\balp + \int_{\mathfrak{P}_l(H)\cap (\mathcal{L}_{l}(L;H))^c}|F(\boldsymbol{\alpha})|^{18}\, d\balp.
\end{aligned}
\end{equation}
% By the symmetry of variables $x$ and $y$, we notice that it suffices to show the cases $l=1$ and $2.$ 
We first consider the case $l=1.$ 
Let us estimate the first term on the right hand side of $(\ref{new dissection})$. By H\"older's inequality, one has
\begin{equation}\label{mean value over the secondary minor arcs}
\begin{aligned} 
I_1 := &\int_{\mathfrak{P}_1(H)\cap \mathcal{L}_{1}(L;H)}|F(\boldsymbol{\alpha})|^{18}\, d\balp \\
    & \le  \left(\int_{\mathfrak{P}_1(H)\cap \mathcal{L}_{1}(L;H)}|F(\boldsymbol{\alpha})|^{20}\, d\balp\right)^{5/6}\left(\int_{\mathfrak{P}_1(H)\cap \mathcal{L}_{1}(L;H)}|F(\boldsymbol{\alpha})|^{8}\, d\balp\right)^{1/6}.
\end{aligned}
\end{equation}
By the definitions of $\mathcal{M}_1(Q)$ and $\mathcal{L}^*_{2,3}(L;H)$, it follows that
\begin{equation}\label{mean value over the secondary minor arcs2}
    \begin{aligned}
     & I_1 \le \left(\int_{\mathcal{M}_1(H/2)\cap  \mathcal{L}_{1}(L;H)}|F(\boldsymbol{\alpha})|^{20}\, d\balp\right)^{5/6}\left(\int_0^1\int_{\mathfrak{M}(H)^2}\int_0^1|F(\boldsymbol{\alpha})|^{8}\, d\balp\right)^{1/6}.
    \end{aligned}
\end{equation}
Now by applying Lemma \ref{improved minor} with $l=1$ and $Q=H/2$, together with (\ref{J10}) and Lemma \ref{k-1 dim measure}, we deduce that
\begin{equation}\label{final bound 1}
    \begin{aligned}
         I_1  &\ll \bigl(X^{8+\varepsilon}J_{10}(2X)\cdot \left(H^{-1}+X^{-1}\right)\cdot L^{-1}\bigr)\bigr)^{5/6}\cdot (X^{10+\varepsilon}\cdot \text{mes}(\mathfrak{M}(H)^2)^{1/6}\\
   & \ll X^{24+\varepsilon}\left(\left(H^{-1}+X^{-1}\right)\cdot L^{-1}\right)^{5/6}\cdot H^{2/3},
    \end{aligned}
\end{equation}
where we have used the fact that $L^{-1}\geq X^{-1/2}$ since $L=H^{1/7}$ and $H\leq X^{3/2}.$

Let us turn to estimating the second term on the right hand side of (\ref{new dissection}). By H\"older's inequality, we have
\begin{equation}\label{mean value over small measure}
\begin{aligned}
    I_2 :=&\int_{\mathfrak{P}_1(H)\cap (\mathcal{L}_{1}(L;H))^c}|F(\boldsymbol{\alpha})|^{18}\, d\balp \\
    & \le  \left(\int_{\mathfrak{P}_1(H)\cap (\mathcal{L}_{1}(L;H))^c}|F(\boldsymbol{\alpha})|^{20}\, d\balp\right)^{5/6}\left(\int_{\mathfrak{P}_1(H)\cap (\mathcal{L}_{1}(L;H))^c}|F(\boldsymbol{\alpha})|^{8}\, d\balp\right)^{1/6}.
\end{aligned}
\end{equation}
Note that 
\begin{equation}\label{P1capL23(L) inclusion}
\begin{aligned}
    \mathfrak{P}_1(H)\cap (\mathcal{L}_{1}(L;H))^c\subseteq  \mathfrak{P}_1(H)\subseteq \mathfrak{m}(H/2)\times [0,1)^3,
\end{aligned}
\end{equation}
by the definition of $\mathfrak{P}_1(H).$ Furthermore,  since \begin{equation*}
    \mathfrak{P}_1(H)\subseteq [0,1)\times \mathfrak{M}(H)\times [0,1)^2
\end{equation*}
and
\begin{equation*}
\begin{aligned}
  &\left([0,1)\times \mathfrak{M}(H)\times [0,1)^2\right)\cap (\mathcal{L}_{1}(L;H))^c \\&=\{\boldsymbol{\alpha}\in [0,1)^4:\ \alpha_2\in \mathfrak{M}(H),\ \alpha_3\in \mathfrak{L}^{(\alpha_2)}(L;H)\},
\end{aligned}
\end{equation*}
we note that
\begin{equation}\label{P1capL23(L)c inclusion}
\begin{aligned}
   &\mathfrak{P}_1(H)\cap (\mathcal{L}_{1}(L;H))^c\\
   &\subseteq \left([0,1)\times \mathfrak{M}(H)\times [0,1)^2\right)\cap (\mathcal{L}_{1}(L;H))^c\\
   &\subseteq [0,1)\times \mathfrak{L}(L;H)\times [0,1),
\end{aligned}
\end{equation}
where $ \mathfrak{L}^{}(L;H)$ is defined in Proposition $\ref{pro5.2}$. Therefore,  it follows by $(\ref{mean value over small measure})$, (\ref{P1capL23(L) inclusion}), and $(\ref{P1capL23(L)c inclusion})$, and  that
\begin{equation}\label{final bound 2}
\begin{aligned}
   I_2  &\leq \left(\int_{\mathfrak{m}(H/2)\times [0,1)^3}|F(\boldsymbol{\alpha})|^{20}\, d\balp \right)^{5/6} \left(\int_0^1\int_{\mathfrak{L}^{}(L;H)}\int_0^1|F(\boldsymbol{\alpha})|^{8}\, d\balp \right)^{1/6}\\
    &\ll ((X^{8+\varepsilon}J_{10}(2X)\cdot\left(H^{-1}+X^{-1}\right))^{5/6}\cdot (X^{10+\varepsilon}\cdot \text{mes}(\mathfrak{L}^{}(L;H)))^{1/6}\\
   & \ll X^{24+\varepsilon}\left(H^{-1}+X^{-1}\right)^{5/6}\cdot (H^{3}L^2)^{1/6},
\end{aligned}
\end{equation}
where we have used Lemma $\ref{basic minor}$, Lemma  \ref{k-1 dim measure}, and Proposition $\ref{pro5.2}.$

Therefore, we conclude from $(\ref{new dissection})$ together with (\ref{final bound 1}) and $(\ref{final bound 2})$ that 
\begin{equation}\label{final bound 6}
\begin{aligned}
&\int_{\mathfrak{P}_1(H)}|F(\boldsymbol{\alpha})|^{18}\, d\balp \ll X^{24+\varepsilon}(\Xi_1+\Xi_2),
\end{aligned}
\end{equation}
where 
$$\Xi_1=\left(\left(H^{-1}+X^{-1}\right)\cdot L^{-1}\right)^{5/6}\cdot H^{2/3}$$
and
$$\Xi_2=\left(H^{-1}+X^{-1}\right)^{5/6}\cdot (H^{3}L^2)^{1/6}.$$
On recalling that $L=H^{1/7}$ and $X^{\delta/100}\leq H\leq X^{3/2},$ one infers that 
\begin{equation}\label{Xibd}
\Xi_i \ll H^{-2/7} + X^{-5/6}H^{23/42} \ll H^{-2/7} + X^{-1/84} \qquad (i=1,2).
\end{equation}
It follows that
\begin{equation}\label{final bound 8}
\int_{\mathfrak{P}_1(H)}|F(\boldsymbol{\alpha})|^{18}\, d\balp \ll X^{24-\gamma},
\end{equation}
for some $\gamma>0$, which completes the case $l=1.$

Next, we consider the case $l=2$. Since the steps are almost the same as in the case $l=1,$ we proceed quickly.  
%We first see that 
%\begin{equation}\label{new dissection2}
%\begin{aligned}
%&\int_{\mathfrak{P}_2(H)}|F(\boldsymbol{\alpha})|^{18}\, d\balp =\int_{\mathfrak{P}_2(H)\cap \mathcal{L}_{2}(L;H)}|F(\boldsymbol{\alpha})|^{18}\, d\balp +\int_{\mathfrak{P}_2(H)\cap (\mathcal{L}_{2}(L;H))^c}|F(\boldsymbol{\alpha})|^{18}\, d\balp .
%\end{aligned}
%\end{equation}
%Let us estimate the first term on the right hand side of $(\ref{new dissection})$. 
By following the argument leading from (\ref{mean value over the secondary minor arcs}) to $(\ref{mean value over the secondary minor arcs2})$, we infer that
\begin{equation*}
\begin{aligned}
 I_3 :=&  \int_{\mathfrak{P}_2(H)\cap \mathcal{L}_{2}(L;H)}|F(\boldsymbol{\alpha})|^{18}\, d\balp \\ &\le \left(\int_{\mathcal{M}_2(H/2)\cap \mathcal{L}_{2}(L;H)}|F(\boldsymbol{\alpha})|^{20}\, d\balp \right)^{5/6} \left(\int_0^1 \int_{\mathfrak{M}(H)^2} \int_0^1 |F(\boldsymbol{\alpha})|^{8}\, d\balp\right)^{1/6}.
  %I(\mathfrak{P}_2(H)\cap \mathcal{L}_{2}(L;H),18) \\
  % &\leq    \left( I(\mathcal{M}_2(H/2)\cap \mathcal{L}_{2}(L;H),20)\right)^{5/6} \cdot\left(I([0,1)\times \mathfrak{M}(H)^2\times[0,1),8) \right)^{1/6}.
\end{aligned}
\end{equation*}
By Lemma \ref{improved minor} with $l=2$ and $Q=H/2$, along with (\ref{J10}) and Lemma \ref{k-1 dim measure}, we find just as in (\ref{final bound 1}) that
\begin{equation}\label{final bound 3}
    \begin{aligned}
  %      I_3   &\ll (X^{8+\varepsilon}J_{10}(2X)\cdot\left(H^{-1}+X^{-1}\right)\bigl(L^{-1}+X^{-1/2}\bigr))^{5/6}\\
  % & \ \ \ \ \ \ \ \ \ \ \ \ \ \ \ \ \ \ \ \ \ \ \ \ \ \ \ \ \ \ \ \ \ \ \ \ \ \ \ \ \ \ \ \ \ \ \ \ \ \ \ \ \cdot (X^{10+\varepsilon}\cdot \text{mes}(\mathfrak{M}(H)\times\mathfrak{M}(H)))^{1/6}\\
   I_3 & \ll X^{24+\varepsilon}\left(\left(H^{-1}+X^{-1}\right)\cdot L^{-1}\right)^{5/6}\cdot H^{2/3}.
    \end{aligned}
\end{equation}

%Let us turn to estimate the second term in the right hand side of (\ref{new dissection}). 
Likewise, by following the argument leading from $(\ref{mean value over small measure})$ to $(\ref{final bound 2})$, one infers that 
\begin{equation*}\label{final bound 4}
    \begin{aligned}
  I_4 :=  &\int_{\mathfrak{P}_2(H)\cap (\mathcal{L}_{2}(L;H))^c}|F(\boldsymbol{\alpha})|^{18}\, d\balp \\
     &\leq \left(\int_{[0,1)\times\mathfrak{m}(H/2)\times [0,1)^2}|F(\boldsymbol{\alpha})|^{20}\, d\balp \right)^{5/6} \left(\int_0^1\int_{\mathfrak{L}(L;H)}\int_0^1|F(\boldsymbol{\alpha})|^{8}\, d\balp \right)^{1/6}.
    \end{aligned}
\end{equation*}
Therefore, by applying Lemma $\ref{basic minor}$, Lemma  \ref{k-1 dim measure}, and Proposition $\ref{pro5.2}$ just as in (\ref{final bound 2}), we obtain
%the bound in (\ref{final bound 4}) becomes
\begin{equation}\label{final bound 5}
    \begin{aligned}
 %      I_4 &\ll ((X^{8+\varepsilon}J_{10}(2X)\cdot\left(H^{-1}+X^{-1}\right))^{5/6}\cdot (X^{10+\varepsilon}\cdot \text{mes}(\mathfrak{L}^{}(L;H)))^{1/6}\\
  I_4 &\ll X^{24+\varepsilon}\left(H^{-1}+X^{-1}\right)^{5/6}\cdot (H^{3}L^2)^{1/6}.
    \end{aligned}
\end{equation}

Thus, by substituting (\ref{final bound 3}) and $(\ref{final bound 5})$ into the right hand side of $(\ref{new dissection})$, one has
\begin{equation*}
\begin{aligned}
&\int_{\mathfrak{P}_2(H)}|F(\boldsymbol{\alpha})|^{18}\, d\balp \ll X^{24+\varepsilon}(\Xi_1+\Xi_2),
\end{aligned}
\end{equation*}
with $\Xi_1$ and $\Xi_2$ as defined in $(\ref{final bound 6}).$
On recalling (\ref{Xibd}), one sees that
%that $L=H^{1/7}$ and $X^{\delta/100}\leq H\leq X^{3/2},$ one sees that 
\begin{equation}\label{final bound 7}
\int_{\mathfrak{P}_2(H)}|F(\boldsymbol{\alpha})|^{18}\, d\boldsymbol{\alpha} \ll X^{24-\gamma},
\end{equation}
for some $\gamma>0$, which completes the case $l=2.$ 

It is easily seen that the verifcation $(\ref{claim})$ for the cases $l=3$ and $l=4$ follows a nearly identical pattern. Hence the proof of the lemma is completed by reference to $(\ref{final bound 8})$ and $(\ref{final bound 7})$.
\end{proof}

\end{document}
